\documentclass[a4paper,11pt]{article}

\usepackage{revy}
\usepackage{amssymb}
\usepackage{amsmath}
\usepackage{commath}
\usepackage{amsfonts}

\revyname{MatematikRevyen}
\revyyear{2025}
\version{0.1}
\eta{? minutter}
\status{Færdig}

\title{Et Knækket Kridt}
\author{Hugo '21}

\begin{document}
\maketitle

\begin{roles}
\role{I}[Sommer] Instruktør
\role{In}[Adam] Interviewer
\role{T2}[Frigg] "Thais", Kendis jakkesæt, sorte solbriller
\role{T1}[Thais] Thais
\role{R}[Anna] Rus
\role{S}[Hugo] Studerende


\end{roles}

\begin{props}
    \prop{1} Fine interview stole.
\end{props}


\begin{sketch}

\scene{En rus sidder i kantinen, med kalkulus og den lille blå. S1 og Thais kommer ind med dummit/foote under armene.}

\says{S1} Hvad så dit knækkede kridt? 

\says{Rus} Øhhh jeg løser bare mine lektier.

\says{Thais} Hah. Det hedder ikke lektier - men ugeopgaver eller sådan noget.

\says{S1} Vi hopper på Cafeen? Skal du med og hakke GT?

\says{Rus} Øhh, skal altså lige nå mine opgaver

\says{S1} De jo nemme?

\scene{Hiver en GT op og rækker den mod RUS}

\says{S1} Bund eller resten i Kalkulus!

\scene{Rus ryster på hovedet og gør sig lille. S1 tager hans kalkulus og hælder øllen i, samtidig med at han på en eller anden måde er klam med Kalkulussen.}

\scene{Lys går ned og publikum er forhåbentligt lidt forvirret. AV'en lyser op og intro til MatematikRevy Behind the Sketch kommer på}
\scene{Noget video spiller (Måske Asger Afbryder. I mellem tiden bliver bord båret ud, og scenen transformeres til studie}

\scene En vært/interviewer går ind på scenen talkshow style, flot men simpelt tøj, publikum klapper og værten takker og leverer de første replikker stående.

\says{In} Godaften og velkommen til "Bag om Sketchen". Byd velkommen til aftenens gæster!

\scene{Vinker tre revyster med ind på scenen. De har også elegant tøj på. S1 og R er tydeligt nu gode venner, holder om hinandens skuldre etc. D sætter sig til rette i interview stole. Benene bliver demonstrativt krydset.}

\says{In} Når en sketch begynder har man et splitsekund til at male en scene - vi ser en Rus. Og så kommer du ind S1.

\says{S1} Jeg siger: "Hva så dit knækkede kridt". Man er ikke i tvivl - jeg er top dawg, En sætning, hvor de andre er lemmaer.

\says{In} Det er jo en sketch om ulige relationer, om at finde sin plads som rus og turde sige fra.

\says{Rus}[Bornholmsk (eller anden umiskendelig dansk dialekt)] At spille en rus er svært. At finde sin plads i en ny verden, hvor alle andre ved det man er i gang med at lære.

\says{In} Og så sker det ingen forventer. Øl i Kalkulus!

\says{S1} Den gule Kalkulus er jo symbolet på Rus og uskyldighed. Og jeg ødelægger det. Gruppepres \act{vifter med hånden}, øl, det korruperer. Edens have er lukket.

\says{In} Og så er der Thais.
\says{Rus} Hvad ville revyen være ud Thais?

\says{S1} Han er jo i sandhed en alsidig revyst. Ninja, Rapper, Skuespiller. Sjov at høre på. Sjov at se på.

%\says{R1} (Lidt prætentiøst) Revyen var i dødvande, der var brug for noget nyt. En ny retning. Og ja, endelig fik vi en sand revystjerne.

%\says{I} Lad os se nogle af de ikoniske øjeblikke.

%\scene{Klip af Stine der synger 48 kroner - zoom ind i hjørnet hvor Thais er ninja og blæser. Cut til andre Thais highlights eg. Hvorfor er der to af dem, Matrevy-sound,... (max 10-15 sekunder)}


\says{In} Hvad gør Thais særlig?

\says{S1} På den ene side er hans øjne er diamanter. På den anden side har han skæg. Han er sin egen modsætning - ninjaen man \textbf{gerne} vil se.\\
Hvis en hvilken som helst anden havde sagt replikken "det er ikke lektier - men ugeopgaver", jamen så ville man, ja, fnise.\\
\act{kunstpause, det et vidunder} Men når Thais siger det med hans kække ansigt, hans latterlige figur... Ja så græder vi alle af grin.

\says{In} Og ja, der er jo ingen Thais uden dig T1.

\says{T1} Ja det er mig der er Thais. Jeg har spillet Thais i revyen i 4 år nu og det har været en fuldstændigt vanvittig rolle.

\says{R}\act{Stadig med sin accent} Jeg er så imponeret over din body transformation.

\says{T1} Ja hver blok 1 tager jeg 10 cm på i højden, bare for at tabe dem igen efter lørdagsforestillingen.

\says{In} Vi kan lige smide et før og efter billede op til seerne.

\scene{Billede på AV af T1 hvor der står før og så et billede af Thais i det samme tøj på efter.}

\says{In} Er det ikke utroligt, man kan jo ikke se forskel? Og så har jeg hørt du method-acter?

\says{T1} Ja, hver dag møder jeg op på studiet som Thais. 
Grænsen mellem mig og Thais bliver mere og mere udtværet. Fordelen er at jeg scorer virkelig mange damer som Thais.

\says{In} Men er det ikke også udmattende?

\says{T1} Jo! Folk vil hele tiden spille Whist med mig - jeg kan ikke engang reglerne. Det hårdeste har nok været at gennemføre bacheloren. Jeg er jo bare uddannet på skuespilskolen.

\says{In} Skuespilskolen! Så du kan ikke engang noget matematik??

\says{T1} Talteori er jo ikke så svært...
\end{sketch}

\end{document}

